\documentclass[10pt,a4paper]{amsart}
\usepackage[margin=19mm]{geometry}
\usepackage{amsmath,amssymb,mathtools}
\usepackage[scr=rsfs,cal=euler]{mathalfa}
\usepackage{xfrac}
\usepackage[unicode,hidelinks,bookmarks=false]{hyperref}
\newcommand{\ie}{i.e.\ }
\newcommand{\cf}{cf.\ }
\newcommand{\resp}{respectively\ }
\newcommand{\Subsection}{\subsection}
\providecommand{\nobreakdash}{\nobreak\hbox{-}}
\setlength{\emergencystretch}{2em}
\multlinegap0pt
\usepackage{bm}
\usepackage{bbm}
\usepackage{dsfont}
\usepackage{xspace}
\usepackage{cancel}
\usepackage{mathtools}
\usepackage{amsthm}
\makeatletter
\@ifundefined{theorem}{\newtheorem{theorem}{Theorem}}{}
\@ifundefined{lemma}{\newtheorem{lemma}[theorem]{Lemma}}{}
\makeatother
\def\JB{\mathsf{JB}}
\def\Lefttop{\mathsf{Lefttop}}
\def\first{\mathsf{first}}
\newcommand{\rmnum}[1]{\romannumeral #1}
\title[Bijective proofs of several conjectures on Jacobi permutatio]{Bijective proofs of several conjectures on Jacobi permutations}
\author{Zhicong Lin, Yongzhou Wen,, Sherry H.F. Yan}
\date{Source: arXiv:2608.11771v1 (CC BY 4.0)}
\begin{document}
\maketitle

\noindent\emph{Source and licence.} Bijective proofs of several conjectures on Jacobi permutations. Version arXiv:2608.11771v1 (CC BY 4.0), distributed under the Creative Commons Attribution 4.0 International licence, as recorded in the arXiv licence metadata for this exact version and shown on the abstract page https://arxiv.org/abs/2608.11771v1. Full rights notice: licenses/arxiv-2608.11771-CC-BY-4.0.txt.\par\smallskip

\noindent\textit{Editorial scope.} Counted unit: the Lemma whose source label is \texttt{lem-phi-reverse} in the article (source0.tex lines 1941--1943, source label \texttt{lem-phi-reverse}), delivered together with the article's own complete original proof of it (source0.tex lines 1944--1992, one \texttt{proof} environment of 49 lines). No mathematics is written, repaired or re-derived: statement and proof are the article's own text line for line. Required context retained uncounted: lem-phi (lines 1912--1923). Every definition, display and label of those lines is retained unchanged. Omitted from this package and not redistributed: 1--1911, 1924--1940, 1993--3209. Nothing inside a retained excerpt is omitted or altered: every retained body is delivered exactly as the article's lines are written, so no omission receipt of any kind accompanies this package. Citations: the retained excerpts carry no citation command at all, so no bibliography entry of the article is transcribed and no reference is redistributed. Reference adaptation: none was needed and none was made; every reference inside the delivered unit resolves inside it, so no unresolved reference or citation is produced. Printed slips kept literal: no printed word, symbol, hypothesis or number of the counted statement or of its proof was altered. Layout and class: the article's own class is \texttt{article} with packages including amsmath, tikz, amssymb, amsfonts, float, booktabs, adjustbox, array, graphicx, multirow, color, amsthm, pgfplots, xcolor; the wrapper is the lane's pinned \texttt{amsart} editorial wrapper at 10pt on A4, which restates the theorem-like environments the retained text needs and adds the article's own macro definitions that the retained text needs (\texttt{\textbackslash JB}, \texttt{\textbackslash Lefttop}, \texttt{\textbackslash first}, \texttt{\textbackslash rmnum}), with extra wrapper packages (bm, bbm, dsfont, xspace, cancel, mathtools). The delivered numbering is the unit's own. Assets: no retained span contains a figure environment, a graphics-inclusion command, a PDF-page inclusion, a table or a TikZ picture; no image file is copied into this package, the package declares no asset dependency and no image file is redistributed with it. Licence: version arXiv:2608.11771v1 of this article is distributed under the Creative Commons Attribution 4.0 International licence, as recorded in the arXiv licence metadata for this exact version and shown on the abstract page https://arxiv.org/abs/2608.11771v1. A full rights notice is delivered at licenses/arxiv-2608.11771-CC-BY-4.0.txt. Characters: every retained excerpt is pure ASCII, so the delivered engine, XeLaTeX with its default Latin Modern text font, reports no missing character.
   \begin{lemma}\label{lem-phi}
   	Given  a Jacobi tree $T\in \JB_S$,  we have \begin{equation}\label{eq-phi}
   		(\Lefttop, \first)\phi(T)=(\Lefttop, \first)T.
   	\end{equation} 
   	Moreover, the resulting tree $\phi(T)$ verifies the following  properties.
   	\begin{itemize}
   		\item[\upshape(\rmnum{1})] The root of the tree $\phi(T)$  is an Andr\'e  $\mathrm{I}$ node.
   		
   		\item[\upshape(\rmnum{2})]  The root of $\phi(T)$ is a Jacobi node if and only if the root of $T$ is an  Andr\'e $\mathrm{I}$ node.
   		\item[\upshape(\rmnum{3})]    Both the left subtree and the right subtree of the root in $\phi(T)$ are Jacobi trees.
   		\end{itemize}
   	\end{lemma}

   	 		\begin{lemma}\label{lem-phi-reverse}
   	 		For any two Jacobi trees $T$ and $T'$,  if $\phi(T)=\phi(T')$, then we have $T=T'$. 
   	 	\end{lemma}

   	 	\begin{proof}
   	 		
   	 		In order to prove the assertion, it  suffices to  show that for any tree $T\in \JB_S$, we can recover $T$ from $\phi(T)$.   
   	 		
   	 		
   	 		If the root of $\phi(T)$ is both   an Andr\'e   $\mathrm{I}$  node and a Jacobi node, then 
   	 		by (\rmnum{2}) of Lemma~\ref{lem-phi},  the root of $T$ is an Andr\'e  $\mathrm{I}$  node.  By the construction of $\phi(T)$, we have $\phi(T)=T$. In this case, we can recover $T$ from $\phi(T)$ by letting $T=\phi(T)$.  
   	 		
   	 		
   	 		Otherwise,  the root of $T$ is not an Andr\'e   $\mathrm{I}$  node.  In this case,   find the topmost node $y$ on the left arm of $T$  such that the node $\max(T)$  belongs to the right subtree of  $y$.      
   	 		Find the node $z_0$ on the right chain starting with the node $y$ such that either $\max(T)=z_0$ or  the node $\max(T)$ belongs to the left subtree of the node $z_0$. Assume that when applying the map $\phi$ to $T$, we generate a sequence  of pairs $(T^{(0)}, z_0), (T^{(1)}, z_1)$, $\ldots, (T^{(k)}, z_k)$ , where $T^{(0)}=T$, $\phi(T)=T^{(k)}$, $(T^{(1)}, z_1)=\alpha(T^{(0)}, z_0)$, and  $(T^{(i+1)}, z_{i+1})=\beta(T^{(i)}, z_i)$ for all $1\leq i<k$.  In order to show that  we can recover $T$ from $\phi(T)$, it suffices to show that we can recover   $T^{(i)}$ from    $T^{(i+1)}$ for all $0\leq i<k$.   
   	 		
   	 		
   	 		From the construction of each $T^{(i+1)}$,  one can easily check that the resulting tree $T^{(i+1)}$ verifies the following desired properties. 
   	 		\begin{itemize}
   	 			\item[\upshape{(\rmnum{1})}]  The node $z_{i+1}$ is the unique  node 
   	 			on the left arm of $T^{(i+1)}$ so that the node $\max(T^{(i+1)})$ belongs to the right subtree of $z_{i+1}$;
   	 			\item[\upshape{(\rmnum{2})}]   Either $\max(T^{(i+1)})=z_i$ or the node $\max(T^{(i+1)})$ belongs to the left subtree of the node $z_i$. 
   	 			\item[\upshape{(\rmnum{3})}]  The node $y$ is the left child of the node $z_{i+1}$.
   	 			\item[\upshape{(\rmnum{4})}] We have $z_i>y$ if $i=0$ and $z_i<y$ otherwise. 
   	 		\end{itemize}
   	 		
   	 		Now we are in the position to give a description of the procedure to recover the
   	 		tree $T^{(i)}$ from $T^{(i+1)}$.  
   	 		
   	 		\begin{itemize}
   	 			\item Find the unique node $u$ on the left arm of $T^{(i+1)}$ so that the node $\max(T^{(i+1)})$ belongs to the right subtree of the node $u$.   Assume that the node $u$ has the left child $v$ in $T^{(i+1)}$.
   	 			\item Find the node $w$ on the right chain starting with the node $u$ so that  either $w=\max(T^{(i+1)})$ or the node $\max(T^{(i+1)})$ belongs to the left subtree of the node $w$.  
   	 			\item Generate a tree by distinguishing the following two cases.
   	 			\begin{itemize}
   	 				\item {\bf Case 1: $w>v$. } 
   	 				\begin{itemize}
   	 					\item Let $R$ be the subtree obtained from $T^{(i+1)}_{w}$ by removing the  right subtree of $w$.
   	 		 Remove the subtree $R$ from $T^{(i+1)}$.
   	 					\item Attach the subtree $R$ to the resulting tree by  inserting the node $w$ onto the right chain starting from the node $v$ so that the labels along this chain increase from left to right.
   	 				\end{itemize}
   	 				
   	 				\item {\bf Case 2: $w<v$.}
   	 				\begin{itemize}
   	 					\item Let $R$ be the  left subtree of the node $w$ in $T^{(i+1)}$.
   	 				 Remove the subtree $R$ from $T^{(i+1)}$.
   	 					\item  Insert the node $w$ onto the left arm of $T$ so that the labels along this chain increase  from top to bottom.
   	 					\item Attach   the subtree $R$ to the node $w$ as a right subtree of the node $w$ in the resulting tree. 
   	 				\end{itemize}
   	 				
   	 			\end{itemize}
   	 		\end{itemize}
   	 		Properties (\rmnum{1})-(\rmnum{4}) of $T^{(i+1)}$ ensure that when applying the above procedure to $T^{(i+1)}$,  we will have $u=z_{i+1}$, $v=y$ and $w=z_i$.   This guarantees that the above procedure reverses the transformation from $T^{(i)}$ to $T^{(i+1)}$. Therefore, applying the above procedure to $T^{(i+1)}$ yields the original tree $T^{(i)}$. This completes the proof.  
   	 		\end{proof}

\end{document}
